\documentclass[10pt,a4paper]{amsart}
\usepackage[margin=19mm]{geometry}
\usepackage{amsmath,amssymb,mathtools}
\usepackage[scr=rsfs,cal=euler]{mathalfa}
\usepackage{xfrac}
\usepackage[unicode,hidelinks,bookmarks=false]{hyperref}
\newcommand{\ie}{i.e.\ }
\newcommand{\cf}{cf.\ }
\newcommand{\resp}{respectively\ }
\newcommand{\Subsection}{\subsection}
\providecommand{\nobreakdash}{\nobreak\hbox{-}}
\setlength{\emergencystretch}{2em}
\multlinegap0pt
\usepackage{bm}
\usepackage{bbm}
\usepackage{dsfont}
\usepackage{xspace}
\usepackage{cancel}
\usepackage{mathtools}
\usepackage{amsthm}
\makeatletter
\@ifundefined{theorem}{\newtheorem{theorem}{Theorem}}{}
\@ifundefined{claim}{\newtheorem{claim}[theorem]{Claim}}{}
\@ifundefined{lemma}{\newtheorem{lemma}[theorem]{Lemma}}{}
\makeatother
\DeclareMathOperator{\Cay}{Cay}
\title[Intervals of hypergraph Tur\'an densities]{Intervals of hypergraph Tur\'an densities}
\author{Xizhi Liu, Oleg Pikhurko}
\date{Source: arXiv:2605.25914v1 (CC BY 4.0)}
\begin{document}
\maketitle

\noindent\emph{Source and licence.} Intervals of hypergraph Tur\'an densities. Version arXiv:2605.25914v1 (CC BY 4.0), distributed under the Creative Commons Attribution 4.0 International licence, as recorded in the arXiv licence metadata for this exact version and shown on the abstract page https://arxiv.org/abs/2605.25914v1. Full rights notice: licenses/arxiv-2605.25914-CC-BY-4.0.txt.\par\smallskip

\noindent\textit{Editorial scope.} Counted unit: the Lemma whose source label is \texttt{lem:centered} in the article (source0.tex lines 419--424, source label \texttt{lem:centered}), delivered together with the article's own complete original proof of it (source0.tex lines 426--520, one \texttt{proof} environment of 95 lines). No mathematics is written, repaired or re-derived: statement and proof are the article's own text line for line. Required context retained uncounted: eq:mixing (lines 313--316). Every definition, display and label of those lines is retained unchanged. Omitted from this package and not redistributed: 1--312, 317--418, 425--425, 521--1342. Nothing inside a retained excerpt is omitted or altered: every retained body is delivered exactly as the article's lines are written, so no omission receipt of any kind accompanies this package. Citations: the retained excerpts carry no citation command at all, so no bibliography entry of the article is transcribed and no reference is redistributed. Reference adaptation: none was needed and none was made; every reference inside the delivered unit resolves inside it, so no unresolved reference or citation is produced. Printed slips kept literal: no printed word, symbol, hypothesis or number of the counted statement or of its proof was altered. Layout and class: the article's own class is \texttt{article} with packages including geometry, amsmath, amssymb, amsthm, mathtools, microtype, parskip, tikz, scalerel, hyperref; the wrapper is the lane's pinned \texttt{amsart} editorial wrapper at 10pt on A4, which restates the theorem-like environments the retained text needs and adds the article's own macro definitions that the retained text needs (\texttt{\textbackslash Cay}), with extra wrapper packages (bm, bbm, dsfont, xspace, cancel, mathtools). The delivered numbering is the unit's own. Assets: no retained span contains a figure environment, a graphics-inclusion command, a PDF-page inclusion, a table or a TikZ picture; no image file is copied into this package, the package declares no asset dependency and no image file is redistributed with it. Licence: version arXiv:2605.25914v1 of this article is distributed under the Creative Commons Attribution 4.0 International licence, as recorded in the arXiv licence metadata for this exact version and shown on the abstract page https://arxiv.org/abs/2605.25914v1. A full rights notice is delivered at licenses/arxiv-2605.25914-CC-BY-4.0.txt. Characters: every retained excerpt is pure ASCII, so the delivered engine, XeLaTeX with its default Latin Modern text font, reports no missing character.
\begin{equation}
        \langle f,Mg\rangle \le Dn\,\bar f\,\bar g+\theta\|f^\circ\|_2\|g^\circ\|_2,
        \label{eq:mixing}
\end{equation}

\begin{lemma}\label{lem:centered}
For every $j$, the $3$-graph $H_j$ satisfies $\lambda(H_j)\le 2/9$. Moreover, for every probability vector $\mathbf{y}$ on $V(H_j)$,
\[
        P_{H_j}(\mathbf{y})\le \frac{2D}{3m_j} +\theta\left(\sum_{v\in V(H_j)}y_v^2-\frac1{m_j}\right).
\]
\end{lemma}

\begin{proof}
The inequality $\lambda(H_j)\le 2/9$ follows straightforwardly from the fact that $H_j$ is $3$-partite. It therefore suffices to prove the second assertion of Lemma~\ref{lem:centered}.

	Fix $j$, and write $\Gamma=\Gamma_j$, $S=S_j$, $n=|\Gamma|$, $m=3n$, and $X=X_j$, $Y=Y_j$, $Z=Z_j$. Let $\mathbf{y}=(y_v)_{v\in V(H_j)}$ be a probability vector. Denote its restrictions to $X,Y,Z$, under our identifications of these sets with $\Gamma$, by $\mathbf{a}=(a_x)_{x\in\Gamma}$, $\mathbf{b}=(b_y)_{y\in\Gamma}$ and $\mathbf{c}=(c_z)_{z\in\Gamma}$, respectively.
Define
\[
        A\coloneqq \sum_{x\in\Gamma}a_x, \qquad B\coloneqq \sum_{y\in\Gamma}b_y, \qquad C\coloneqq \sum_{z\in\Gamma}c_z,
\]
so that $A+B+C=1$. Also put
\[
        R_X\coloneqq \sum_{x\in\Gamma} a_x^2, \qquad R_Y\coloneqq \sum_{y\in\Gamma} b_y^2, \qquad R_Z\coloneqq \sum_{z\in\Gamma} c_z^2,
\]
and
\[
        \Delta_X\coloneqq R_X-\frac{A^2}{n}, \qquad \Delta_Y\coloneqq R_Y-\frac{B^2}{n}, \qquad \Delta_Z\coloneqq R_Z-\frac{C^2}{n}.
\]
The quantities $\Delta_X,\Delta_Y,\Delta_Z$ are non-negative by Cauchy's inequality.
Define 
\[
        T
        \coloneqq \sum_{x,y\in\Gamma}\sum_{s\in S}a_xb_yc_{xsy} .
\]
Note that $P_{H_j}(\mathbf{y})=6T$.

We use the standard terminology that the \emph{link} of a vertex $v$ in a $3$-graph is the graph whose edges are the pairs that form hyperedges together with $v$. For a bipartite graph with parts $U$ and $V$, its biadjacency matrix is the $U\times V$ matrix whose $(u,v)$-entry is $1$ if $uv$ is an edge, and is $0$ otherwise. 

\begin{claim}\label{clm:link-relabel}
For every vertex of $H_j$, its link between the other two vertex classes is a bipartite graph whose biadjacency matrix, after relabellings that are permutations of $\Gamma$, is identical to the adjacency matrix of $\Cay(\Gamma,S)$.
\end{claim}

\begin{proof}[Proof of the claim]
Fix $y_\ast\in Y$. The corresponding link has parts $X$ and $Z$. A pair $(x,z)$ is an edge precisely when $z=xsy_\ast$ for some $s\in S$. If we relabel the $X$-side by $\hat{x}\coloneqq x$ and the $Z$-side by $\hat{z}\coloneqq zy_\ast^{-1}$, then this condition becomes $\hat{z}=\hat{x}s$.

Fix $x_\ast\in X$. The corresponding link has parts $Y$ and $Z$. Relabel the $Y$-side by $\hat{y}\coloneqq y^{-1}$ and the $Z$-side by $\hat{z}\coloneqq z^{-1}x_\ast$. The relation $z=x_\ast sy$ is equivalent to $\hat{z}=\hat{y}s^{-1}$, which has the same right Cayley form since $S=S^{-1}$.

Finally, fix $z_\ast\in Z$. The corresponding link has parts $Y$ and $X$. Relabel the $Y$-side by $\hat{y}\coloneqq y^{-1}$ and the $X$-side by $\hat{x}\coloneqq z_\ast^{-1}x$. The relation $z_\ast=xsy$ is equivalent to $\hat{x}=\hat{y}s^{-1}$, again of the same right Cayley form. 

In each case the relabellings are permutations of $\Gamma$, and the resulting biadjacency relation is exactly the right Cayley relation.
\end{proof}

For fixed $y\in\Gamma$, define the vector $\mathbf{c}^{(y)}=(c^{(y)}_v)_{v\in\Gamma}$ by $c^{(y)}_v\coloneqq c_{vy}$. By this definition,
\[
        \sum_{x\in\Gamma}\sum_{s\in S}a_xc_{xsy} = \sum_{x\in\Gamma}\sum_{s\in S}a_xc^{(y)}_{xs}.
\]
The right-hand side is $\langle \mathbf{a},M\mathbf{c}^{(y)}\rangle$, where $M=M_{\Gamma,S}$ is the adjacency matrix of $\Cay(\Gamma,S)$. We view $\mathbf{a}$ and $\mathbf{c}^{(y)}$ as functions on $\Gamma$. The vector $\mathbf{a}$ has total mass $A$ and centred $L^2$ norm $\sqrt{\Delta_X}$. Also, since $v\mapsto vy$ is a permutation of $\Gamma$, the vector $\mathbf{c}^{(y)}$ has total mass $C$ and centred $L^2$ norm $\sqrt{\Delta_Z}$. Hence \eqref{eq:mixing} gives
\[
        \sum_{x\in\Gamma}\sum_{s\in S}a_xc_{xsy}
        = \langle \mathbf{a},M\mathbf{c}^{(y)}\rangle
        \le \frac{D}{n}AC+\theta\sqrt{\Delta_X\Delta_Z}.
\]
Multiplying this inequality by $b_y$ and summing over $y\in\Gamma$, using $\sum_{y\in\Gamma} b_y=B$, gives
\begin{equation}
        T\le \frac{D}{n}ABC+\theta B\sqrt{\Delta_X\Delta_Z}.
        \label{eq:XZ}
\end{equation}
The same argument applied to the links obtained by fixing a vertex in $X$ and in $Z$, using the relabellings from Claim~\ref{clm:link-relabel}, gives the analogous bounds
\begin{equation}
        T\le \frac{D}{n}ABC+\theta A\sqrt{\Delta_Y\Delta_Z},
        \label{eq:YZ}
\end{equation}
and
\begin{equation}
        T\le \frac{D}{n}ABC+\theta C\sqrt{\Delta_X\Delta_Y}.
        \label{eq:XY}
\end{equation}
Averaging \eqref{eq:XZ}, \eqref{eq:YZ} and \eqref{eq:XY}, and using the inequality $2\sqrt{uv}\le u+v$, we obtain
\begin{align*}
        T
        &\le \frac{D}{n}ABC
        +\frac{\theta}{3}
        \left(
        B\sqrt{\Delta_X\Delta_Z}
        +A\sqrt{\Delta_Y\Delta_Z}
        +C\sqrt{\Delta_X\Delta_Y}
        \right) 
        \le \frac{D}{n}ABC
        +\frac{\theta}{6}(\Delta_X+\Delta_Y+\Delta_Z).
\end{align*}
Multiplying by $6$, we get
\begin{equation}
        P_{H_j}(\mathbf{y})
        = 6T
        \le \frac{6D}{n}ABC +\theta(\Delta_X+\Delta_Y+\Delta_Z).
        \label{eq:lambdaH-main}
\end{equation}
Now $6ABC\le 2/9$, because $A+B+C=1$, and
\[
        \Delta_X+\Delta_Y+\Delta_Z =\sum_{v\in V(H_j)}y_v^2-\frac{A^2+B^2+C^2}{n} \le \sum_{v\in V(H_j)}y_v^2-\frac1{3n}.
\]
Since $m=3n$, \eqref{eq:lambdaH-main} implies
\[
        P_{H_j}(\mathbf{y}) \le \frac{2D}{9n} +\theta\left(\sum_{v\in V(H_j)}y_v^2-\frac1{3n}\right) = \frac{2D}{3m} +\theta\left(\sum_{v\in V(H_j)}y_v^2-\frac1m\right).
\]
This proves the centered estimate.
\end{proof}

\end{document}
